\section*{Introduction}
\phantomsection\addcontentsline{toc}{section}{Introduction}

Before building a complete grading system, two questions must be answered. First, does
the GraphSAGE template with structural descriptors behave as a general architectural
property, or does it only work on one particular dataset? Second, do global structural and
topological features of the retina actually carry information that a deep network trained
on the raw image does not already capture? If either answer were negative, the systems of
the following chapters would rest on weak foundations.

This chapter answers both questions with the first contributions of the thesis. In
Part~A, we validate the graph template on distributed-systems data, where graphs are
natural and ground truth is unambiguous: microservice orchestration and anomaly detection
on six benchmarks. We then explain how each element of this template is transferred to
retinal grading. In Part~B, we move to diabetic retinopathy and test the value of
structural features with three independent methods: a hybrid CNN--GNN model with
topological features, a gradient-boosting model on hand-crafted structural features
(SAMF-GB), and a reduced CPU-efficient variant (LightGBM-HC). Weaker intermediate results
are reported deliberately, to show how the final solution was built step by step.

\section{Part A: Validating the Graph Template Outside Medicine}
\label{sec:ch2-rationale}

Part~A validates the GraphSAGE architectural template on six structurally diverse,
non-medical benchmarks before any contact with retinal images. The rationale is
epistemological: if the structural-invariant augmentation works across six domains with
very different graph topologies, the expectation that it will transfer to inter-patient
similarity graphs rests on reproducible evidence rather than on analogy alone. This
six-benchmark validation is Contribution~C1; it also makes the architectural choices of
TOPORET and DOTS easier to defend before a regulatory reviewer.

\paragraph{Why distributed systems?}
Validating a graph architecture requires data where the graph is natural and the labels
are unambiguous. Microservice architectures and computer networks provide exactly this:
the dependencies between services or hosts form an explicit graph, and failures or attacks
are recorded without the inter-grader disagreement that affects medical labels. They also
share an essential property with a screening service: the system must make decisions on
new nodes as they appear, without being retrained. A template that works in these
conditions is a credible candidate for a population graph of patients, and any weakness
found there can be corrected before touching medical data.

\subsection{GNN-MSOrchest: Microservice Cascade-Failure Recovery}
\label{sec:msorchest}

GNN-MSOrchest~\cite{belhadj2025icaart} models a 48-service e-commerce microservice system
as a directed attributed graph $G=(V,E,X)$. Each node $v\in V$ carries a five-dimensional
health vector $x_v=[\mathrm{CPU},\mathrm{mem},\mathrm{lat},\mathrm{err},\mathrm{dep}]^\top$
updated every $200$~ms. Edges encode directed dependencies (service $u$ calls service
$v$) with propagation-delay attributes. A two-layer GraphSAGE with mean aggregation drives
adaptive workload allocation:
\begin{equation}
\hat y_v = \mathrm{MLP}\bigl(h^{(2)}_v\bigr),
\label{eq:msorch}
\end{equation}
where $h^{(2)}_v$ is the node representation after two message-passing rounds with
$d_h=256$ and $S=25$ sampled neighbours per layer. The prediction $\hat y_v$ of \cref{eq:msorch} drives
load-shedding decisions for node $v$ during cascade propagation.

\paragraph{Structural-invariant augmentation.}
The key methodological contribution is the augmentation
$h_{\mathrm{aug}}=[\,h_{\mathrm{GNN}}\,\Vert\,I(v)\,]$, where $I(v)$ is a
four-dimensional structural descriptor:
\begin{equation}
I(v) = \bigl(\delta(v),\, c(v),\, b(v),\, h(v)\bigr)^\top ,
\label{eq:invariants}
\end{equation}
with $\delta(v)$ the degree centrality, $c(v)$ the clustering coefficient, $b(v)$ the
betweenness centrality and $h(v)$ the local structural entropy. These quantities depend
only on graph topology, not on node features or learned representations, and are cheap
to compute (linear in $|V|+|E|$ for all but betweenness, which is approximated by
sampling).

\Cref{tab:msorch} compares GNN-MSOrchest with a static-threshold baseline. The cascade
isolation mechanism -- withholding load from nodes whose predictive entropy exceeds a
threshold -- is the direct methodological precursor of the entropy abstention of DOTS
(\cref{sec:abst}).

\begin{table}[H]
\centering
\caption[GNN-MSOrchest versus a static-threshold baseline]{GNN-MSOrchest versus a static-threshold baseline (SimPy simulation, 48
nodes)~\cite{belhadj2025icaart}.}
\label{tab:msorch}
\small
\begin{tabular}{@{}lcc@{}}
\toprule
Metric & Static-threshold baseline & GNN-MSOrchest \\
\midrule
Mean response time (ms) & 142 & 98 ($-31.0\%$) \\
p99 latency (ms) & 510 & 287 ($-43.7\%$) \\
Cascade-failure containment & 34\% & 81\% ($+47$~pp) \\
Resource-utilisation Gini coefficient & 0.41 & 0.18 ($-56.1\%$) \\
\bottomrule
\end{tabular}
\end{table}

\subsection{Structural Invariants Across Six Benchmarks}
\label{sec:invariants}

The second publication~\cite{belhadj2026invariants} tests whether the augmentation of
\cref{eq:invariants} gives statistically significant and stable gains across six
benchmarks with heterogeneous graph structure: a synthetic controlled graph, two
network-intrusion datasets (KDD Cup~1999~\cite{tavallaee2009nslkdd},
CICIDS~2017~\cite{sharafaldin2018cicids}), a more recent network-intrusion
dataset (UNSW-NB15~\cite{moustafa2015unsw}), a refined intrusion benchmark
(NSL-KDD~\cite{tavallaee2009nslkdd}) and a production
microservice trace (SockShop/Istio).

\Cref{tab:invariants} reports the six results. Each row is one benchmark: its domain,
the F1 score of the plain GNN, the F1 score once the four invariants $I(v)$ are appended
to the node features, and the resulting gain in recall on the anomalous class. The recall
gain is positive on every benchmark and lies between $+3.1$ and $+5.0$~pp (mean $+4.2$~pp, standard deviation $0.6$~pp). The
coefficient of variation ($0.6/4.2\approx 14\%$) is low, so the effect is stable across
domains rather than driven by a single benchmark.

\begin{table}[H]
\centering
\caption[Structural-invariant augmentation across six graph benchmarks]{Structural-invariant augmentation across six graph
benchmarks~\cite{belhadj2026invariants}.}
\label{tab:invariants}
\small
\begin{tabular}{@{}llccc@{}}
\toprule
Benchmark & Domain & F1 (GNN) & F1 ($+I(v)$) & $\Delta$ recall \\
\midrule
Synthetic graph & Controlled network & 85.3\% & 93.5\% & $+5.0$~pp \\
KDD Cup 1999 & Network intrusion (classic) & 88.1\% & 94.2\% & $+4.6$~pp \\
CICIDS 2017 & Multi-protocol IDS & 86.7\% & 92.0\% & $+4.1$~pp \\
UNSW-NB15 & Network intrusion (recent) & 84.9\% & 90.8\% & $+3.1$~pp \\
NSL-KDD & Intrusion (refined) & 87.4\% & 93.1\% & $+4.4$~pp \\
SockShop/Istio & Production microservices & 89.2\% & 95.3\% & $+4.4$~pp \\
\midrule
Mean (SD) & & 87.0\% (1.5) & 93.2\% (1.6) & $+4.2$~pp (0.6) \\
\bottomrule
\end{tabular}
\end{table}

\subsection{Transfer Architecture and Its Justification}
\label{sec:transfer}

The transfer from distributed systems to DR grading is direct at the level of the
architecture: GraphSAGE with $L=2$, $d_h=256$, dropout $p=0.30$, Adam with learning rate
$10^{-3}$ and $S=25$ is used unchanged. The structural augmentation
$h_{\mathrm{aug}}=[\,h_{\mathrm{GNN}}\Vert I(v)\,]$ becomes
$z_i=[\,x_i^{\mathrm{CNN}}\Vert x_i^{\mathrm{TDA}}\,]$, and the cascade isolation of
uncertain nodes becomes entropy abstention at threshold $\tau_U$.

\paragraph{From distributed systems to retinal grading.}
The transfer rests on a one-to-one correspondence between the objects of the two
domains, set out in \cref{tab:mapping} in pipeline order. The first three rows define the
nodes and what they carry: a service becomes a patient, the service health vector
becomes the CNN embedding of the fundus image, and the structural invariant $I(v)$
becomes the topological feature vector of \cref{eq:tda}. The fourth row defines the
edges: a dependency between services becomes a dual-similarity edge between patients
(\cref{eq:dualsim}). The last three rows define behaviour: a cascade failure becomes
diagnostic uncertainty spreading across a grade boundary, load shedding becomes
abstention, and recall-prioritised anomaly detection becomes the class-weighted CORAL
loss that favours the rare severe grades. Read row by row, the table is a checklist:
every component validated on network graphs has a named counterpart in TOPORET
(\cref{ch:toporet}) and DOTS (\cref{ch:dots}).

\begin{table}[H]
\centering
\caption[Mapping from distributed-systems concepts to DR-grading analogues]{Conceptual mapping from distributed-systems concepts (this chapter) to their
DR-grading analogues (\cref{ch:toporet,ch:dots}).}
\label{tab:mapping}
\small
\begin{tabular}{@{}P{6.4cm}P{8.0cm}@{}}
\toprule
Distributed-systems concept & DR-grading analogue \\
\midrule
Service node $v\in V$ & Patient $i$ \\
Service health vector (CPU, memory, latency, \ldots) & CNN embedding $x_i^{\mathrm{CNN}}$ \\
Structural invariant $I(v)$ (degree, clustering, betweenness, entropy) & Topological feature vector $x_i^{\mathrm{TDA}}$ (\cref{eq:tda}) \\
Dependency edge (service $u$ calls service $v$) & Population-graph edge (dual similarity $S_{ij}$, \cref{eq:dualsim}) \\
Cascade failure (overload propagating across services) & Diagnostic uncertainty propagating across a grade boundary \\
Load-shedding decision (withhold traffic from an overloaded node) & Abstention (defer to a specialist when $H_i>\tau_U$) \\
Recall-prioritised anomaly detection (rare-event sensitivity) & Class-weighted CORAL loss ($\lambda_3=8$, $\lambda_4=6$) \\
\bottomrule
\end{tabular}
\end{table}

\paragraph{Why hyperparameters transfer without re-tuning.}
A natural objection is that hyperparameters validated on network-traffic and microservice
graphs may not suit a population graph of patients. Two considerations mitigate this.
First, the graphs of Part~A already span a wide range of sizes (from the 48-node SimPy
testbed of \cref{tab:msorch} to the traffic graphs of \cref{tab:invariants}, with several
thousand nodes) and densities, and the augmentation gain of \cref{tab:invariants} is stable
across the six benchmarks. Second, the population
graphs of \cref{ch:toporet,ch:dots} (several thousand to tens of thousands of nodes,
$k=15$ neighbours per node) fall within that range, so the transfer is an interpolation
rather than an extrapolation. The hyperparameter sensitivity analysis was nevertheless
repeated on DR data (\cref{tab:hparam}) as a confirmatory check.

\paragraph{What the transfer claims and does not claim.}
The transfer is architectural, not distributional. It does not claim that retinal
population graphs and microservice graphs share statistical properties, nor that
performance on one guarantees performance on the other. It claims that the same
architectural pattern -- GraphSAGE with mean aggregation, $L=2$ hops, $d_h=256$,
structural-invariant augmentation -- transfers without modification, which is evidence
for the genericity of the approach. The cross-domain validation is methodological
insurance, not a proof of universal applicability.

\subsection{Threats to the Transfer Argument}
\label{sec:threats}

Three threats deserve to be named. First, all six benchmarks are graphs over discrete
events or entities (network flows, service calls); a population graph over patients
carries node features of a different statistical character, and stability across event
graphs does not by itself guarantee stability across patient-similarity graphs. Second,
the benchmarks share an implicit assumption that better-connected neighbourhoods are more
informative. Its analogue for disease severity -- that similar patients share diagnostic
information -- is precisely what H1 tests, not a premise this chapter can establish.
Third, six benchmarks remain a finite, non-random sample of graph domains; the
coefficient-of-variation argument supports stability within this sample only.

These caveats do not undermine the strategy; they are why the hyperparameter analysis is
repeated on retinal data (\cref{sec:dualsim}) instead of relying on Part~A alone.

\section{Part B: Structural Evidence on Retinal Data}
\label{sec:partb}

Part~B presents three convergent proof-of-concept validations of the structural axis
(the answer to L2). Three methodologies -- a deep hybrid model, gradient boosting with SHAP attribution, and
a leave-one-family-out ablation -- each indicate that global structural priors improve DR
grading.
This convergence motivates the formal statistical test of H2 in \cref{ch:toporet}.

\paragraph{Experimental setting.}
The three studies of Part~B use the public datasets of \cref{sec:datasets}. The hybrid
CNN--GNN model was developed on APTOS~2019 with a binary target (No DR versus any DR),
which allowed many ablation runs at low cost. SAMF-GB and LightGBM-HC were evaluated on the
five-grade task on Kaggle EyePACS, the benchmark used for TOPORET and DOTS in the following
chapters, with five-fold cross-validation stratified by grade. SAMF-GB and LightGBM-HC
extract the vessel network with the same pre-processing (CLAHE and Zhang--Suen
skeletonisation); the hybrid model instead computes its descriptors from
structure-enhanced intensity responses, without explicit vessel segmentation. The three
studies therefore differ in data, task and input representation, and their results are
compared qualitatively rather than numerically.

\subsection{Hybrid CNN--GNN with Geometric and Topological Invariants}
\label{sec:hybrid}

The first retinal study~\cite{belhadj2026hybrid}, extended in the book
chapter~\cite{mezghich2026lnaitopo}, tests whether global structural descriptors add
information to a deep CNN--GNN pipeline. An EfficientNet-B0 backbone extracts local visual
features; a graph neural network with mean aggregation models the relations between the
positions of the CNN feature map, connected by a four-neighbourhood rule; and two families of
global descriptors are added: curvature-based geometric descriptors of the vessels and
persistent-homology invariants ($H_0$ and $H_1$ lifespans) computed from structure-enhanced
intensity responses, without explicit vessel segmentation. The four streams are concatenated
and classified by a fully connected layer.

\paragraph{Protocol.}
The main evaluation uses APTOS~2019 ($3{,}662$ images, resized to $128\times128$ pixels to keep
the persistent-homology computation tractable), with a stratified $70/15/15$ split, on the
clinically relevant binary screening task (No DR versus any DR). Five-class grading is a
secondary analysis, and Messidor-2 ($1{,}748$ images) is used, without fine-tuning, to test
generalisation under domain shift.

\paragraph{Ablation.}
\Cref{tab:hybrid} adds each component in turn to the CNN baseline. Relational modelling by
the GNN gives the largest single gain ($+3.4$~pp); geometric and topological descriptors add
$+1.2$ and $+1.7$~pp respectively. The full model reaches $95.1\%$ accuracy (AUC $0.970$,
95\% bootstrap interval $[94.0\%,\,96.2\%]$), $+7.0$~pp over the baseline, which is more than
the sum of the individual gains ($6.3$~pp): the components interact synergistically, and
McNemar's test confirms the improvement over the CNN baseline ($p<0.01$). On Messidor-2,
without any fine-tuning, the full model reaches $92.1\%$ accuracy against $88.1\%$ for the
CNN alone, and it keeps a higher AUC under Gaussian noise, contrast changes, blur and
downsampling~\cite{mezghich2026lnaitopo}.

\begin{table}[!tb]
\centering
\caption[Ablation of the hybrid CNN--GNN pipeline]{Ablation of the hybrid CNN--GNN pipeline
on APTOS~2019 (binary screening, held-out test set)~\cite{mezghich2026lnaitopo}.}
\label{tab:hybrid}
\small
\begin{tabular}{@{}lcc@{}}
\toprule
Configuration & Accuracy & $\Delta$ vs.\ CNN \\
\midrule
CNN only (EfficientNet-B0) & 88.1\% & -- \\
CNN $+$ geometric descriptors & 89.3\% & $+1.2$~pp \\
CNN $+$ topological descriptors & 89.8\% & $+1.7$~pp \\
CNN $+$ GNN & 91.5\% & $+3.4$~pp \\
CNN $+$ GNN $+$ geometry and topology (full) & \textbf{95.1\%} & $\mathbf{+7.0}$~pp \\
\bottomrule
\end{tabular}
\tabnote{Sum of the individual gains: $1.2+1.7+3.4=6.3$~pp, below the $7.0$~pp of the full
model.}
\end{table}

\paragraph{Five-class grading and training behaviour.}
On the harder five-class task, the full model reaches $87.5\%$ accuracy, with macro
precision, recall and F1 all at $87.3\%$. \Cref{fig:hybridcm} shows where the errors lie:
No~DR, Mild and Moderate are well separated, and most errors are confusions between the two
advanced grades, Severe and Proliferative, which share vascular pathology at the coarse
$128\times128$ resolution. \Cref{fig:hybridcurves} shows the training loss and accuracy:
the loss decreases steadily and the accuracy saturates in the last epochs, without the
divergence that would indicate unstable optimisation.

\begin{figure}[!tb]
\centering
\includegraphics[width=0.66\linewidth]{hybrid_confusion_5class}
\caption[Confusion matrix of the hybrid model on five-class grading]{Confusion matrix of the
hybrid CNN--GNN model on five-class grading (APTOS~2019). Correct predictions lie on the
diagonal; most errors are between Severe and Proliferative DR. Adapted
from~\cite{mezghich2026lnaitopo}.}
\label{fig:hybridcm}
\end{figure}

\begin{figure}[!tb]
\centering
\includegraphics[width=0.78\linewidth]{hybrid_training_curves}
\caption[Training curves of the hybrid model]{Training loss and training accuracy of the
hybrid CNN--GNN model on five-class grading, per epoch. Adapted
from~\cite{mezghich2026lnaitopo}.}
\label{fig:hybridcurves}
\end{figure}

\paragraph{Topological descriptors of the vessel skeleton.}
The later systems of this thesis describe the vascular network with a fixed vector of
persistence statistics. For each fundus image, the vessel skeleton is extracted with CLAHE
pre-processing~\cite{zuiderveld1994clahe} followed by Zhang--Suen
thinning~\cite{zhang1984thinning}. The branch points of the skeleton form a point cloud
$P_i=\{p_1,\ldots,p_n\}\subset\mathbb{R}^2$. The Vietoris--Rips filtration is computed
with GUDHI~\cite{maria2014gudhi}, which gives a seven-dimensional feature vector
\begin{equation}
x_i^{\mathrm{TDA}} = \bigl(\Pi_{H_0},\, n_{H_0},\, \Pi_{H_1},\, n_{H_1},\, \beta_1,\,
D_f,\, \Delta\bigr)^\top \in \mathbb{R}^7 ,
\label{eq:tda}
\end{equation}
where $\Pi_{H_k}$ is the total $k$-th persistence, $n_{H_k}$ the number of components
or loops, $\beta_1$ the first Betti number, $D_f$ the fractal dimension, and $\Delta$ the
maximum $H_1$ persistence (a neovascularisation signal).

The synergy observed in \cref{tab:hybrid} motivates the next step. In the hybrid model,
the graph links positions within one image and the topological descriptors are simply
concatenated; TOPORET (\cref{sec:dualsim}) instead builds a graph between patients and uses
the distance between persistence diagrams inside the similarity that defines its edges,
fusing both signals at graph-construction time.

\subsection{SAMF-GB: Structural Features Outperform GPU EfficientNet-B3}
\label{sec:samfgb}

SAMF-GB~\cite{belhadj2026samfgb} shows that seven families of hand-crafted structural
features, combined with LightGBM gradient boosting~\cite{ke2017lightgbm}, outperform a
GPU EfficientNet-B3 at about $1/13$ of the computational cost. The 256-dimensional
feature vector is composed as follows:
\begin{itemize}
\item \textbf{vascular morphology} (48-D): vessel calibre reduction and tortuosity;
\item \textbf{fractal dimension} (12-D): loss of vascular self-similarity;
\item \textbf{persistent homology} (7-D): topology of the vessel network (\cref{eq:tda});
\item \textbf{DWT texture} (64-D, four levels): microaneurysm and exudate texture;
\item \textbf{Gabor responses} (32-D): vessel edges and IRMA characteristics;
\item \textbf{radial FFT} (24-D): degradation of the spoke-wheel vessel pattern;
\item \textbf{LBP $+$ log-Euclidean} (69-D): haemorrhage and exudate cluster patterns.
\end{itemize}

\Cref{tab:shap} gives the dimension and mean absolute SHAP value~\cite{lundberg2017shap}
of each family. The three most discriminative families are all global structural
features (vascular morphology, fractal dimension, persistent homology), not local
appearance features. This is independent evidence for H2, obtained with a method entirely
different from the ANOVA of \cref{sec:h2}.

\begin{table}[H]
\centering
\caption[SAMF-GB feature families and mean absolute SHAP values]{SAMF-GB feature families, dimension and mean absolute SHAP
value~\cite{belhadj2026samfgb}.}
\label{tab:shap}
\small
\begin{tabular}{@{}lccP{6.2cm}@{}}
\toprule
Family & Dim. & Mean $|\mathrm{SHAP}|$ & Primary DR correlate \\
\midrule
Vascular morphology & 48 & 0.142 & Vessel calibre reduction, tortuosity \\
Fractal dimension & 12 & 0.138 & Loss of vascular self-similarity \\
Persistent homology & 7 & 0.131 & Topology of the vascular network \\
DWT (4 levels) & 64 & 0.094 & Microaneurysm and exudate texture \\
Gabor responses & 32 & 0.087 & Vessel edges and IRMA \\
Radial FFT & 24 & 0.082 & Spoke-wheel vessel pattern degradation \\
LBP $+$ log-Euclidean & 69 & 0.068 & Haemorrhage and exudate clusters \\
\midrule
Full SAMF-GB & 256 & -- & QWK $=0.882$; $2.1$~ms CPU inference \\
\bottomrule
\end{tabular}
\end{table}

The SHAP ranking of \cref{tab:shap} is also a check on the clinical plausibility of the
model: the three families that matter most describe the geometry and topology of the
vessels, which are the structures that ophthalmologists examine when they look for venous
beading, capillary closure and new vessels.

\Cref{tab:samf} compares SAMF-GB with
EfficientNet-B3: SAMF-GB gains $+0.042$ QWK, its classifier runs on a 16-core CPU instead
of a GPU at about 13 times lower cost, and its predictions can be explained with SHAP.

\begin{table}[H]
\centering
\caption[SAMF-GB versus EfficientNet-B3 on Kaggle EyePACS]{SAMF-GB versus EfficientNet-B3 on Kaggle EyePACS ($N=88{,}702$).}
\label{tab:samf}
\small
\begin{tabular}{@{}lcccl@{}}
\toprule
System & QWK & Inference & Hardware & Art.~13 support \\
\midrule
EfficientNet-B3 & 0.840 & 85~ms & NVIDIA A100 & No (opaque) \\
SAMF-GB & 0.882 & 2.1~ms & 16-core CPU & Yes (SHAP per prediction) \\
\midrule
Difference & $+0.042$ & $\times 40$ (classifier only) & $\times 13$ cheaper & -- \\
\bottomrule
\end{tabular}
\tabnote{The $2.1$~ms of SAMF-GB is the inference time of the gradient-boosting model on
pre-computed features. Feature extraction is not included: the persistent-homology stage
alone takes about $18$~ms per image on the same CPU (\cref{tab:latency}). The speed ratio
therefore compares classifiers, not complete pipelines.}
\end{table}

\subsection{LightGBM-HC}
\label{sec:lgbmhc}

Publication~\cite{belhadj2026icpram} validates a reduced 128-dimensional hand-crafted
feature set with LightGBM, reaching QWK $=0.858$ at $1.6$~ms CPU inference on the same
benchmark. The hyperparameters, each selected by five-fold cross-validated grid search,
are listed in \cref{tab:lgbm-hp}. The QWK is lower than that of SAMF-GB ($0.858$ versus $0.882$, a difference of $0.024$).
LightGBM-HC omits the 7-D persistent-homology features, but also the radial FFT and
LBP families and most of the Gabor responses; the difference therefore cannot be
attributed to topology alone. It is consistent with, but does not isolate, a contribution
of the topological component; the isolated contribution of the TDA features is measured
in the TOPORET ablation ($+0.016$ QWK, \cref{tab:ablation3}).

\begin{table}[H]
\centering
\caption[LightGBM-HC hyperparameters]{LightGBM-HC hyperparameters, selected by five-fold cross-validated grid search on
Kaggle EyePACS (search range in parentheses).}
\label{tab:lgbm-hp}
\small
\begin{tabular}{@{}ll@{}}
\toprule
Hyperparameter & Selected value (search range) \\
\midrule
Number of boosting rounds & 480 (100--1000) \\
Maximum tree depth & 7 (3--12) \\
Learning rate & 0.04 (0.01--0.2, log scale) \\
Number of leaves & 63 (15--127) \\
Minimum child samples & 20 (5--50) \\
Feature subsampling fraction & 0.75 (0.5--1.0) \\
L2 regularisation & 1.2 (0--5) \\
\bottomrule
\end{tabular}
\end{table}

\paragraph{Feature-family ablation.}
To isolate the contribution of each of the four families retained in the 128-D set
(vascular morphology, fractal dimension, DWT texture, Gabor responses), a
leave-one-family-out ablation was run (\cref{tab:lgbm-abl}). Vascular morphology is by
far the most load-bearing family ($-0.077$ QWK when removed), more than the two
texture-oriented families together. This agrees with the SHAP ranking of SAMF-GB and
reaches the same conclusion by a different method: global vascular structure carries more
signal than local texture for this task.

\begin{table}[!tb]
\centering
\caption[Leave-one-family-out ablation for LightGBM-HC]{Leave-one-family-out ablation for LightGBM-HC. Each row removes one family from
the 128-D set.}
\label{tab:lgbm-abl}
\small
\begin{tabular}{@{}lcc@{}}
\toprule
Configuration & QWK & $\Delta$ vs.\ full set \\
\midrule
Full 128-D feature set & 0.858 & -- \\
$-$ Vascular morphology (48-D) & 0.781 & $-0.077$ \\
$-$ Fractal dimension (12-D) & 0.829 & $-0.029$ \\
$-$ DWT texture (64-D) & 0.842 & $-0.016$ \\
$-$ Gabor responses (4-D, reduced) & 0.851 & $-0.007$ \\
\bottomrule
\end{tabular}
\end{table}

With only four families and $128$ dimensions, LightGBM-HC keeps most of the accuracy of
SAMF-GB at a lower cost, which makes it a practical option for screening sites equipped only
with standard computers. Both gradient-boosting models also provide feature-level explanations directly, without any post-hoc method.

\subsection{Convergent Evidence from Four Independent Methodologies}
\label{sec:converge}

Four methodologically independent approaches --- the three studies of Part~B and the
formal ANOVA of \cref{ch:toporet} --- reach the same qualitative conclusion:
global structural priors improve DR grading.
\begin{enumerate}
\item \textbf{Deep learning}~\cite{belhadj2026hybrid,mezghich2026lnaitopo}: topological descriptors add $+1.7$~pp
beyond the CNN baseline (\cref{tab:hybrid}).
\item \textbf{Gradient boosting}~\cite{belhadj2026samfgb}: SAMF-GB outperforms GPU
EfficientNet-B3 by $+0.042$ QWK at about 13 times lower cost; topology ranks third of
seven families by SHAP (\cref{tab:shap}).
\item \textbf{Reduced gradient boosting}~\cite{belhadj2026icpram}: the reduced feature set
without topology (and without three other families) loses $0.024$ QWK, and the
leave-one-family-out ablation ranks global vascular morphology first
(\cref{tab:lgbm-abl}).
\item \textbf{Formal ANOVA} (\cref{sec:h2}): $F(4,4995)=1308.9$, $\eta^2=0.51$.
\end{enumerate}

\paragraph{Why four methodologies, not one.}
One could ask whether a single, larger experiment would have established the same point
with less apparent duplication. Four methodologies were used because each has its own
characteristic failure mode, and agreement across methods with different failure modes is
harder to explain away than one large effect. A deep-learning ablation
(\cref{tab:hybrid}) can be confounded by optimisation idiosyncrasies; a SHAP attribution
(\cref{tab:shap}) by correlated features inflating or deflating attributions; a
leave-one-family-out ablation (\cref{tab:lgbm-abl}) by interactions between the removed
and retained families. No single confound threatens all four at once, so their agreement
is correspondingly hard to attribute to a methodological artefact. The same logic
motivates the distributed-systems detour of Part~A: breadth of independent evidence,
rather than depth within one experimental paradigm, is the organising principle of the
evidential strategy of this thesis. The same principle guides the evaluation of TOPORET and DOTS in the next chapters.

\section{Synthesis of the First Contributions}
\label{sec:ch3-synthesis}

\Cref{tab:ch3summary} summarises the five studies of this chapter.

Three design decisions for the rest of the thesis follow directly from these studies. The
graph template, its hyperparameters and its inductive mode can be reused unchanged on
retinal data, since they proved stable across six domains. Topological descriptors of the
vessel skeleton deserve a place in the model, since three different methods found them
informative. Finally, the way the two signals are combined matters: since they are already synergistic
when simply concatenated, the next system goes further and fuses topology into the
similarity that builds a graph between patients. The remainder of this section summarises the evidence on which these decisions rest.

\begin{table}[H]
\centering
\caption{Summary of the first contributions.}
\label{tab:ch3summary}
\small
\begin{tabular}{@{}P{3.6cm} P{3.4cm} P{5.2cm} c@{}}
\toprule
Study & Data & Main result & Publication \\
\midrule
GNN-MSOrchest & 48-service microservice simulation & Cascade containment $81\%$ vs.\ $34\%$ & P1 \\
Structural invariants & Six graph benchmarks & Mean recall $+4.2$~pp, CV $14\%$ & P2 \\
Hybrid CNN--GNN & APTOS 2019 (binary) & $95.1\%$ accuracy ($+7.0$~pp), AUC $0.970$; synergistic gains & P3, B2 \\
SAMF-GB & Kaggle EyePACS & QWK $0.882$ vs.\ $0.840$; CPU classifier ($2.1$~ms) & J2 \\
LightGBM-HC & Kaggle EyePACS & QWK $0.858$ with a reduced 128-D set & P4 \\
\bottomrule
\end{tabular}
\end{table}

\section*{Conclusion}
\phantomsection\addcontentsline{toc}{section}{Conclusion}

This chapter established three empirical findings. (F1)~The GraphSAGE template with
structural descriptors gives stable gains on six heterogeneous benchmarks
(mean recall gain $+4.2$~pp, coefficient of variation $14\%$). (F2)~On DR data, relational and
topological contributions are complementary: the full hybrid model gains $+7.0$~pp over the
CNN baseline, more than the sum of the individual gains ($6.3$~pp). (F3)~Global
structural and topological features carry discriminative information, as shown
independently by a deep hybrid model, by gradient boosting with SHAP attribution and by a
leave-one-family-out ablation. As a practical by-product, SAMF-GB reaches a higher QWK than a GPU EfficientNet-B3 with a
gradient-boosting classifier that runs on a standard CPU ($2.1$~ms per image once the
features are extracted).

These results do not prove that a population graph will work on retinal images, but they
show that it is worth testing, with topology fused into the similarity that builds the
graph. The next chapter implements this idea in TOPORET.